Vision--language models (VLMs) learn rich image--text associations, yet they reason with poorly localized visual attention: the saliency maps induced by their attention are noisy, place mass on irrelevant regions, and only loosely track the objects a token describes. We introduce the \emph{Saliency Alignment Loss}, a training-time auxiliary objective that supervises a VLM's token-conditioned attention with region-level grounding annotations. For each generated token we read out a saliency map over the image patches and minimize its divergence from a target distribution supported on the token's annotated region, jointly with the standard captioning loss. The objective adds no parameters and leaves the architecture untouched, dropping into ordinary fine-tuning pipelines. Fine-tuning LLaVA-1.5-7B on the densely grounded COCONut-PanCap captions for only $200$ steps, we measure alignment intrinsically on held-out images with three metrics---chance-normalized Attention Mass Ratio, pixel-level Average Precision, and Normalized Scanpath Saliency---chosen to sidestep the COCO contamination that pervades public grounding benchmarks. Alignment transforms saliency---mean Attention Mass Ratio rises $4.6{\times}$ and the maps visibly concentrate on the relevant objects---while leaving language quality intact. We trace \emph{where} this adaptation occurs: the language model, not the multimodal projector, absorbs the signal, reaching a reproducible low-rank solution in a small set of grounding-relevant attention heads. We are equally explicit about its limits: the intrinsic gains do not yet transfer to downstream benchmarks, saliency quality and task grounding are dissociated, and alignment withdraws attention from background ``stuff'' the annotations never cover. These findings motivate an in-progress instruction-tuned variant designed to convert sharper attention into better grounding.
